%%%PAGE 219 rot=0 legibility=good summary=伏安法内外接选择（比较法/试触法）；灯泡各类图像（P-U、P-U²、R-U、P-I、P-I²、P-R、P-1/R）；电阻箱误差；力学实验当成力学计算题；下半页空白
%%%BLOCK id=219-01 ch=10 sec=伏安法·内外接选择 kind=method exp=1 alt=none cont=none y=0.04-0.165
\begin{itemize}
\item \unclear{…}：$R_x\ \rule[0.5ex]{2.5em}{0.4pt}\ \sqrt{R_AR_V}$ \quad 大内偏大 \quad 小外偏小。
% CHECK: 本行行首页边小字（约两字）无法辨认
\item \unclear{试触}：变化大、相差大、接一块，$\dfrac{\Delta U}{U_0}\ \rule[0.5ex]{1.5em}{0.4pt}\ \dfrac{\Delta I}{I_0}$

待测电阻和变化大的表接一块
\end{itemize}
%%%END
%%%BLOCK id=219-02 ch=10 sec=小灯泡伏安特性·图像 kind=diagram exp=1 alt=none cont=none y=0.165-0.49
灯泡

%FIG p219_a 0.03 0.182 0.84 0.49
\notefig[0.85]{figures/p219_a.png}
% CHECK: 第三行左图纵轴标注形似 R（可能应为 P）

（图右上侧文字）\unclear{导数函数} $y=f(x)x$，
%%%END
%%%BLOCK id=219-03 ch=10 sec=电学实验·电阻箱误差 kind=note exp=1 alt=none cont=none y=0.495-0.55
电阻箱可能无法调到 4.nn$\,\Omega$，造成误差：电阻箱精确值\struck{\unclear{过}大}\red{不够}
% CHECK: "4.nn" 原稿为字母 nn，疑为 4.00 之类；红笔在"过(?)大"处划线并在其后以红笔写"不够"，疑为改正
%%%END
%%%BLOCK id=219-04 ch=18 sec=实验方法·力学实验 kind=note exp=1 alt=none cont=none y=0.565-0.61
力学实验当成力学计算题
%%%END
%%%PAGEEND 219
